\documentclass[10pt,a4paper]{amsart}
\usepackage[margin=19mm]{geometry}
\usepackage{amsmath,amssymb,mathtools}
\usepackage[scr=rsfs,cal=euler]{mathalfa}
\usepackage{xfrac}
\usepackage[unicode,hidelinks,bookmarks=false]{hyperref}
\newcommand{\ie}{i.e.\ }
\newcommand{\cf}{cf.\ }
\newcommand{\resp}{respectively\ }
\newcommand{\Subsection}{\subsection}
\providecommand{\nobreakdash}{\nobreak\hbox{-}}
\setlength{\emergencystretch}{2em}
\multlinegap0pt
\usepackage{amssymb}
\usepackage{tikz}
\usepackage{algorithm}\usepackage{algpseudocode}
\newtheorem{theorem}{Theorem}
\title{Numerical Approaches for Identifying the Time-Dependent Potential Coefficient in the Diffusion Equation}
\author{Arshyn Altybay, Michael Ruzhansky}
\date{Source: arXiv:2511.08302v1 (CC BY 4.0)}
\begin{document}
\maketitle

\noindent\emph{Source and licence.} Numerical Approaches for Identifying the Time-Dependent Potential Coefficient in the Diffusion Equation. Version arXiv:2511.08302v1 (CC BY 4.0), distributed by arXiv under the Creative Commons Attribution 4.0 International licence, http://creativecommons.org/licenses/by/4.0/, as recorded in the arXiv licence metadata for this exact version and shown on the abstract page https://arxiv.org/abs/2511.08302v1. Full licence notice: \texttt{licenses/arxiv-2511.08302-CC-BY-4.0.txt}.\par\smallskip

\noindent\textit{Editorial scope.} Counted unit: the article's theorem thm:direct-wellposed (source0.tex lines 208--210). The article's complete original proof of it is delivered with it (source0.tex lines 212--328, one \texttt{proof} environment of 117 lines). No mathematics is written, repaired or re-derived: the statement and the proof are the article's own lines, in the article's own notation, and no retained line is substituted or re-flowed.  Retained uncounted as the source-local context the proof needs: lines 83--85; lines 91--93.  Nothing was omitted from any retained span, so the package carries no omission record.  No retained line contains a citation, so the wrapper carries no bibliography and no bibliography entry.  No retained line contains a figure environment, an image-inclusion command or drawing code, so no image is delivered and the package declares no asset dependency.  Editorial formatting only, in addition: an AMS \texttt{amsart} preamble that restates the article's own declarations and macros used by the retained lines, namely the environments \texttt{theorem} on separate counters, together with the standard packages those lines need.  The retained environments print in this document as Theorem 1 in order of appearance, whereas the article prints the same items with the numbering of its own full document.  The complete native source of the article as distributed in the arXiv e-print for this exact version is bundled unchanged as \texttt{sources/188879/source0.tex}.
\begin{equation} \label{1.1}
    u_t(x,t) = u_{xx}(x,t) - p(t)u(x,t) + f(x,t), \quad (x,t) \in Q := \Omega \times (0,T],
\end{equation}

\begin{equation} \label{1.3}
    u(0,t) = u(l,t) = 0, \quad t \in [0,T].
\end{equation}

\begin{theorem}[Existence and Uniqueness]\label{thm:direct-wellposed}
Let the assumptions (1)-(3) be satisfied. Then the initial-boundary value problem \eqref{1.1}--\eqref{1.3} admits a unique solution $u \in C ([0,T]; H^2(0,l) \cap H_0^1(0,l))$ with $u_t \in L^2(0,T; L^2(0,l))$.
\end{theorem}

\begin{proof}
\textbf{Existence:} The solution is constructed via Fourier series. Consider the eigenfunctions $\phi_n(x) = \sqrt{\frac{2}{l}} \sin\left(\frac{n\pi x}{l}\right)$ and eigenvalues $\lambda_n = \left(\frac{n\pi}{l}\right)^2$ for $n \in \mathbb{N}$. Expand:
\[
u(x,t) = \sum_{n=1}^{\infty} u_n(t) \phi_n(x), \quad
\varphi(x) = \sum_{n=1}^{\infty} \varphi_n \phi_n(x), \quad
f(x,t) = \sum_{n=1}^{\infty} f_n(t) \phi_n(x),
\]
where $\varphi_n = \langle \varphi, \phi_n \rangle$, $f_n(t) = \langle f(\cdot,t), \phi_n \rangle$. Substitution into the PDE yields:
\[
u_n'(t) + (\lambda_n + p(t))u_n(t) = f_n(t), \quad u_n(0) = \varphi_n.
\]
Solving this ODE gives:
\begin{equation}\label{sol_coeff}
u_n(t) = \varphi_n e^{-\lambda_nt-\Lambda(t)} + \int_{0}^{t} e^{-\lambda_n(t-s)-(\Lambda(t)-\Lambda(s))}\, f_n(s)\, ds, \quad  \Lambda(t) =  \int_0^t p(s)ds.
\end{equation}

Thus, the full solution is:
\begin{equation}\label{full}
u(t,x) =\sum_{n=1
}^{\infty}\left[ \varphi_n e^{-\lambda_nt-\Lambda(t)} + \int_{0}^{t} e^{-\lambda_n(t-s)-(\Lambda(t)-\Lambda(s))}\, f_n(s)\, ds\right]\phi_n(x).
\end{equation}

Since $p\in L^\infty(0,T)$, let $P:=\|p\|_{L^\infty(0,T)}$. Then, for $0\le s\le t\le T$,
\[
|\Lambda(t)| = \left|\int_0^t p(s)ds \right| \le Pt,\qquad |\Lambda(t)-\Lambda(s)|\le P(t-s),
\]
and hence
\[
\big|e^{-\Lambda(t)}\big|\le e^{Pt},\qquad \big|e^{-(\Lambda(t)-\Lambda(s))}\big|\le e^{P(t-s)}.
\]

The smoothness and boundary vanishing of $\varphi$ and $f(\cdot,t)$ imply the Fourier–sine
coefficient decay
\[
|\varphi_n|=O(n^{-2}),\qquad \sup_{t\in[0,T]}|f_n(t)|=O(n^{-2}).
\]
Consequently, with $\lambda_n\asymp n^2$, there exist $C, C>0$ such that, for all $t\in[0,T]$,
\[
|u_n(t)|\;\le\; C\!\left(\frac{e^{-c n^2 t}}{n^{2}}+\frac{1}{n^{4}}\right),
\qquad
\lambda_n|u_n(t)|\;\le\; C\!\left(e^{-c n^2 t}+\frac{1}{n^{2}}\right).
\]
In particular, for every $\delta\in(0,T]$,
\[
\sup_{t\in[\delta,T]}|u_n(t)|=O(n^{-4}),\qquad
\sup_{t\in[\delta,T]}\lambda_n|u_n(t)|=O(n^{-2}).
\]

Therefore, for each $\delta>0$, the Fourier series defining $u$, $u_x$, and $u_{xx}$
converge \emph{uniformly} on $[0,l]\times[\delta,T]$.
At $t=0$ one has $u(\cdot,0)=\varphi$ in $L^2(0,l)$.
Moreover,
\[
u_t=u_{xx}-p(t)u+f\in L^2\!\big(0,T;L^2(0,l)\big).
\]


\textbf{Uniqueness:} 
Suppose $u_1$ and $u_2$ are two solutions satisfying \eqref{1.1}--\eqref{1.3}. Define $w(x,t) = u_1(x,t) - u_2(x,t)$. Then $w$ satisfies:
\begin{align*}
    w_t &= w_{xx} - p(t) w, \quad (x,t) \in Q, \\
    w(x,0) &= 0, \quad x \in [0,l], \\
    w(0,t) &= w(l,t) = 0, \quad t \in [0,T].
\end{align*}

Consider the energy functional:
\[
E(t) = \frac{1}{2} \int_0^l |w|^2  dx.
\]
Differentiate with respect to $t$ and substitute the PDE:
\begin{align*}
\frac{d}{dt}E(t) &= Re\int_0^l \bar{w}w_t dx \\
&= Re\int_0^l \bar{w} [w_{xx} - p(t) w]  dx \\
&= Re\int_0^l \bar{w} w_{xx}  dx - p(t) \int_0^l |w|^2  dx.
\end{align*}
Integration by parts (using boundary conditions $w(0,t) = w(l,t) = 0$) yields:
\[
\int_0^l \bar{w} w_{xx}  dx = \left[ \bar{w} w_x \right]_0^l - \int_0^l \bar{w_x} w_x   dx = - \int_0^l |w_x|^2 dx.
\]
So,
\[
\frac{d}{dt}E(t) = - \int_0^l |w_x|^2 dx - p(t) \int_0^l |w|^2 dx
\]
Since $p \in L^{\infty}([0,T])$ there exists $P>0$ such that $|p(t)| \leq P\ \text{for almost every  $t \in [0,T]$,}$ Then,
\[
- p(t) \int_0^l |w|^2 dx \leq P \int_0^l |w|^2 dx
\]
Hence,
\[
\frac{d}{dt}E(t) \leq - \int_0^l |w_x|^2 dx + P \int_0^l |w|^2 dx.
\]
Apply Poincaré's inequality for $H_0^1(0,l)$:
\[
\int_0^l |w|^2 dx  \leq \frac{l^2}{\pi^2} \int_0^l | w_x |^2 dx \implies \int_0^l | w_x |^2 dx \geq \frac{\pi^2}{l^2} \int_0^l|w|^2 dx.
\]
Substitute into the inequality:
\[
\frac{d}{dt}E(t) \leq - \frac{\pi^2}{l^2}\int_0^l |w|^2 dx + P \int_0^l |w|^2 dx = (P-\frac{\pi^2}{l^2})\int_0^l |w|^2 dx.
\]
Since $\int_0^l |w|^2 dx = 2E(t)$, we get:
\[
\frac{d}{dt}E(t) \leq 2(P-\frac{\pi^2}{l^2})E(t).
\]
Let $M=2(P-\frac{\pi^2}{l^2})$. Then
\[
\frac{d}{dt}E(t) \leq ME(t).
\]
Since $ E(0) = 0$ and $E(t) \geq 0$, Gronwall’s inequality implies:
\[
E(t) \leq 0, \text{ for all }  t \in [0,T].
\]
Therefore, $E(t)=0$ for all t, so $w(x,t)=0$ almost everywhere. Hence, 
\[
u_1(x,t) = u_2(x,t),
\]
which completes the proof.
\end{proof}

\end{document}
